% Source presentation for register row N-c-011.

\subsection{Schild's \(\mathcal{ALC}\)--modal correspondence}
\label{sec:p09-presentation-n-c-011}

\subsubsection{The two source languages}

Fix sets \(\mathsf C\) of atomic concept symbols and \(\mathsf R\) of atomic
role symbols.  Schild's concept descriptions have the grammar
\[
 \begin{aligned}
 C,D &::= c\mid\mathsf{Top}\mid C\cap D\mid\neg C\mid\forall R.C,\\
 R &::= r,
 \qquad c\in\mathsf C,\quad r\in\mathsf R.
 \end{aligned}
\]
The source prints the separate role-expression production because Section~4
later enlarges it.  In the present fragment its only production is \(R::=r\),
so the remainder eliminates that one-step category and writes \(\forall r.C\).
This is an eliminable handbook normalization of the displayed source grammar,
not an identification of the later compound-role language with this fragment.
Union and existential restriction are metalinguistic abbreviations, not new
primitive constructors:
\[
 C\cup D := \neg(\neg C\cap\neg D),
 \qquad
 \exists r.C := \neg\forall r.\neg C.
\]
An interpretation \(I\) has a domain \(\Delta^I\), assigns each
\(c\in\mathsf C\) a subset \(c^I\subseteq\Delta^I\), and assigns each
\(r\in\mathsf R\) a binary relation
\(r^I\subseteq\Delta^I\times\Delta^I\).  Here the handbook's
\(I\), \(\Delta^I\), and superscript \((\cdot)^I\) replace Schild's domain
\(D\) and extension function \(\mathcal E\): \(D\) becomes \(\Delta^I\), and
\(\mathcal E[X]\) becomes \(X^I\).  This eliminable normalization separates
the interpretation, its carrier, and its typed extension map and avoids
reusing the source's \(\mathcal E\) below; it changes no semantic clause.
Schild literally permits
\(\Delta^I\) to be any set, including the empty set.  The extension clauses
are
\[
\begin{aligned}
 \mathsf{Top}^I&=\Delta^I,
 & (C\cap D)^I&=C^I\cap D^I,\\
 (\neg C)^I&=\Delta^I\setminus C^I,
 & (\forall r.C)^I
   &=\{x\in\Delta^I:\forall y\in\Delta^I
       ((x,y)\in r^I\Rightarrow y\in C^I)\}.
\end{aligned}
\]
A concept is \emph{coherent} when its extension is nonempty in at least one
interpretation, and is \emph{incoherent} otherwise.  The relation
\(C\sqsubseteq D\) holds when
\(C^I\subseteq D^I\) for every interpretation \(I\), and
\(C\equiv D\) holds when those extensions are equal for every \(I\)
\cite[p.~467]{Schild1991}.  An empty interpretation cannot witness coherence,
and all concept extensions in it are empty; hence retaining Schild's literal
convention changes neither coherence nor the universal inclusion and equality
tests.

Suppose that \(\mathsf R=\{r_1,\ldots,r_m\}\).  The matching normal
multimodal language \(\mathbf K_{(m)}\) has propositional letters \(p_c\),
one box \(K_r\) for each \(r\in\mathsf R\), and grammar
\[
 \alpha,\beta ::= p_c\mid\top\mid\alpha\land\beta
               \mid\neg\alpha\mid K_r\alpha.
\]
Disjunction, implication, and possibility are metalinguistic abbreviations:
\[
 \alpha\lor\beta:=\neg(\neg\alpha\land\neg\beta),
 \qquad
 \alpha\mathbin{\to}\beta:=\neg(\alpha\land\neg\beta),
 \qquad
 \Diamond_r\alpha:=\neg K_r\neg\alpha.
\]
A model is
\(M=(W,(R_r)_{r\in\mathsf R},V)\), where \(W\) is any set, each
\(R_r\subseteq W\times W\), and \(V(p_c)\subseteq W\).  Writing
\(\llbracket\alpha\rrbracket_M\) for the set of worlds satisfying
\(\alpha\), its clauses are
\[
\begin{aligned}
 \llbracket p_c\rrbracket_M&=V(p_c),
 &\llbracket\top\rrbracket_M&=W,\\
 \llbracket\alpha\land\beta\rrbracket_M
   &=\llbracket\alpha\rrbracket_M\cap\llbracket\beta\rrbracket_M,
 &\llbracket\neg\alpha\rrbracket_M
   &=W\setminus\llbracket\alpha\rrbracket_M,\\
 \llbracket K_r\alpha\rrbracket_M
   &=\{x\in W:\forall y\in W
       ((x,y)\in R_r\Rightarrow
        y\in\llbracket\alpha\rrbracket_M)\}.
\end{aligned}
\]
For \(x\in W\), define
\(M,x\models\alpha\) to mean
\(x\in\llbracket\alpha\rrbracket_M\).  A formula \(\alpha\) is
\emph{satisfiable} when
\(\llbracket\alpha\rrbracket_M\ne\varnothing\) for at least one model \(M\),
and is \emph{valid} when \(\neg\alpha\) is not satisfiable.  Equivalently,
\(\alpha\) is valid exactly when every world in every model satisfies it.
Thus an empty world set witnesses neither satisfiability nor invalidity.
Write \(\models_{\mathbf K_{(m)}}\alpha\) exactly when \(\alpha\) is valid in
this just-defined model class.  This is a metalinguistic abbreviation for
validity, not another semantic relation, and is used below to state the
subsumption correspondence.
No frame condition is imposed: this is the normal logic \(\mathbf K\) with
\(m\) modalities, rather than a reflexive, transitive, functional, or otherwise
restricted frame class.

\subsubsection{The translation and what it preserves}

The map \(f\) is a metalinguistic map from the set of
\(\mathcal{ALC}\)-concept descriptions over \(\mathsf C,\mathsf R\) to the
set of \(\mathbf K_{(m)}\)-formulas over the corresponding letters and boxes:
\[
\begin{aligned}
 f(c)&=p_c,&
 f(\mathsf{Top})&=\top,\\
 f(C\cap D)&=f(C)\land f(D),&
 f(\neg C)&=\neg f(C),\\
 f(\forall r.C)&=K_r f(C).&&
\end{aligned}
\]
Schild prints \(f(c)=c\), overloading one atomic name across the two syntactic
categories.  The displayed \(f(c)=p_c\) is the handbook's type-separating
notational normalization: \(p_c\) is the propositional letter corresponding
to the concept symbol \(c\), not an additional semantic object.  Conversely,
replacing each \(p_c\) by \(c\), each \(K_r\) by
\(\forall r\), and the Boolean constructors by their concept counterparts
recovers a unique concept description.  Thus the map is an isomorphism of the
displayed grammars, not merely an embedding into a richer modal language.

Schild says that the preservation claim could be shown by induction on concept
complexity but does not print the induction.  The next argument is the
handbook's reconstruction of that claim.  Given an interpretation \(I\), form
\[
 M_I=
 \bigl(\Delta^I,(r^I)_{r\in\mathsf R},V_I\bigr),
 \qquad V_I(p_c)=c^I.
\]
Conversely, for a modal model
\(M=(W,(R_r)_{r\in\mathsf R},V)\), define the interpretation \(I_M\) by
\[
 \Delta^{I_M}=W,
 \qquad
 c^{I_M}=V(p_c),
 \qquad
 r^{I_M}=R_r.
\]
The latter is a semantic construction, distinct from the inverse syntactic
replacement described above, and \(M_{I_M}=M\) on the displayed vocabulary.
Then, for every concept description \(C\) and every
\(x\in\Delta^I\),
\[
 x\in C^I
 \quad\Longleftrightarrow\quad
 x\in\llbracket f(C)\rrbracket_{M_I}
 \quad\Longleftrightarrow\quad
 M_I,x\models f(C).
\]
The proof is structural induction on \(C\).  The atomic and
\(\mathsf{Top}\) cases are the definitions of \(V_I\) and the common domain.
Intersection and complement commute with \(f\) because both semantics use the
same set-theoretic intersection and relative complement.  For the only modal
case, the induction hypothesis gives
\[
\begin{aligned}
 x\in(\forall r.C)^I
 &\Longleftrightarrow
   \forall y\bigl((x,y)\in r^I\Rightarrow y\in C^I\bigr)\\
 &\Longleftrightarrow
   \forall y\bigl((x,y)\in R_r\Rightarrow
                    M_I,y\models f(C)\bigr)\\
 &\Longleftrightarrow M_I,x\models K_r f(C).
\end{aligned}
\]
Applying the same induction to \(I_M\) gives, for every \(x\in W\),
\[
 x\in C^{I_M}
 \quad\Longleftrightarrow\quad
 M,x\models f(C).
\]
Consequently, \(C\) is coherent exactly when \(f(C)\) is satisfiable.
Because each constructor becomes one constructor, both syntactic directions
increase length by at most a linear factor.  This reconstructs the claim from
which Schild concludes Theorem~1: \(\mathcal{ALC}\) is a notational variant of
\(\mathbf K_{(m)}\), and satisfiability in the latter has the same
computational complexity as coherence in the former
\cite[thm.~1, p.~467]{Schild1991}.

The displayed structural-induction equivalence also yields, as a handbook
consequence rather than an extra clause of Theorem~1,
\[
 C\sqsubseteq D
 \quad\Longleftrightarrow\quad
 \models_{\mathbf K_{(m)}} f(C)\rightarrow f(D).
\]
Indeed, a failure of inclusion in an interpretation becomes a world falsifying
the implication in its associated model, and a modal counterexample becomes a
failure of inclusion in its associated interpretation.  This is the precise
subsumption/validity relationship used below.

\subsubsection{The same expression in three notations}

Before stating the correspondence, Schild illustrates the translations with
the concept
\[
 \neg c_1\cup\forall r.(c_2\cap c_3).
\]
Its first-order rendering in the source is
\[
 \neg c_1(x)\lor
 \forall y\bigl(r(x,y)\mathbin{\Rightarrow}
                  (c_2(y)\land c_3(y))\bigr),
\]
and, using the handbook's type-separating propositional letters, its modal
rendering is
\[
 \neg p_{c_1}\lor K_r(p_{c_2}\land p_{c_3}).
\]
These are three notations for the same recursively matched structure; the
first-order display is explanatory and does not add first-order quantifiers to
the concept grammar \cite[p.~467]{Schild1991}.

\subsubsection{The source's concept example}

Schild's example is not merely the informal paraphrase of a woman whose
children are sons.  Using the source's lowercase plural symbols, it requires a
woman who is not a man, has at least one male child, and has only children that
are men and not women:
\[
 \mathsf{women}\cap\neg\mathsf{men}
 \cap\exists\mathsf{child}.\mathsf{men}
 \cap\forall\mathsf{child}.
      (\mathsf{men}\cap\neg\mathsf{women}).
\]
The source prints that concept expression but not the following modal formula.
Applying \(f\) gives the handbook reconstruction
\[
 p_{\mathsf{women}}\land\neg p_{\mathsf{men}}
 \land\Diamond_{\mathsf{child}}p_{\mathsf{men}}
 \land K_{\mathsf{child}}
      (p_{\mathsf{men}}\land\neg p_{\mathsf{women}}).
\]
The concept is printed on p.~467 of the source \cite[p.~467]{Schild1991}.
The following two-world evaluation is likewise a handbook reconstruction.  Let
\(W=\{w,v\}\) and
\(R_{\mathsf{child}}=\{(w,v)\}\).  The formula holds at \(w\) when
\(w\in V(p_{\mathsf{women}})\setminus V(p_{\mathsf{men}})\) and
\(v\in V(p_{\mathsf{men}})\setminus V(p_{\mathsf{women}})\).  Enlarging the
model with a further \(\mathsf{child}\)-successor of \(w\) outside
\(V(p_{\mathsf{men}})\), or inside \(V(p_{\mathsf{women}})\), falsifies the
universal conjunct in both presentations.  Replacing
\(R_{\mathsf{child}}\) by the empty relation falsifies the existential conjunct
in both presentations.

\subsubsection{The Boolean reductions}

Schild's Lemma~1 applies to any terminological logic comprising Boolean
operations on concepts, not only to the displayed \(\mathcal{ALC}\) language.
It records logspace reductions among subsumption, equivalence, and incoherence.
The source's three reductions can be read directly from the Boolean
constructors:
\[
\begin{aligned}
 C\sqsubseteq D
 &\Longleftrightarrow C\cap D\equiv C,\\
 C\equiv D
 &\Longleftrightarrow
   (C\cap\neg D)\cup(D\cap\neg C)
   \text{ is incoherent},\\
 C\text{ is incoherent}
 &\Longleftrightarrow C\sqsubseteq\neg\mathsf{Top}.
\end{aligned}
\]
Each transformation copies its inputs only a constant number of times and is
therefore computable in logarithmic workspace
\cite[lem.~1, p.~467]{Schild1991}.

\subsubsection{Two immediate modal-logic consequences}

By referral to standard modal-logic model results, Schild first observes that
one may consider either finite \(\mathcal{ALC}\)-models of exponential size or
one infinite canonical \(\mathcal{ALC}\)-model.  The article does not prove the
underlying modal result locally \cite[p.~468]{Schild1991}.

Proposition~1 then records an equational consequence.  Schild prints the
equations and their soundness-and-completeness claim but gives no separate name
to their syntactic generator set.  Reusing his \(\mathcal E\), the semantic
extension function above, would conflate two different kinds of object, while
the semantic relation \(\equiv\) cannot itself name syntactic derivability.
Define the metalinguistic set \(\mathsf{EqAx}_{\mathcal{ALC}}\) to be exactly
the following set of ordered pairs of concept descriptions: first, every pair
\((\sigma(s),\sigma(t))\), where \(s=t\) is an equation in the signature
\((\mathsf{Top},\cap,\neg)\) valid under every valuation in every Boolean
algebra and \(\sigma\) is the homomorphic extension of an assignment of its
variables to concept descriptions; second, for every atomic role \(r\), the
pair
\[
 (\forall r.\mathsf{Top},\mathsf{Top});
\]
and third, for every atomic role \(r\) and concept descriptions \(C,D\), the
pair
\[
 \bigl(\forall r.(C\cap D),
       (\forall r.C)\cap(\forall r.D)\bigr).
\]
Define \(C=_{\mathsf{EqAx}_{\mathcal{ALC}}}D\) when \((C,D)\) belongs to the
least equivalence relation on concept descriptions that contains every ordered
pair in \(\mathsf{EqAx}_{\mathcal{ALC}}\) and is compatible with \(\cap\),
\(\neg\), and every operation \(\forall r\).  Thus replacement in every
concept context is included.  The set
\(\mathsf{EqAx}_{\mathcal{ALC}}\) and the generated relation
\(=_{\mathsf{EqAx}_{\mathcal{ALC}}}\) are metalinguistic, not concept
constructors or semantic relations.  Unlike \(\equiv\), which compares
extensions in every interpretation, the former is a set of syntactic equation
generators and the latter is their least constructor-compatible equivalence
relation.  Both
are needed immediately.  With this handbook reconstruction,
Proposition~1's soundness-and-completeness claim becomes
\[
 C=_{\mathsf{EqAx}_{\mathcal{ALC}}}D
 \quad\Longleftrightarrow\quad C\equiv D.
\]
Thus the stated Boolean-algebra identities and two value-restriction schemes
are sound and complete for \(\mathcal{ALC}\)-equivalence.  Schild refers to
Lemmon's 1966 modal-algebra result; no proof is printed in the article
\cite[prop.~1, p.~468]{Schild1991}.

\subsubsection{Why subsumption is PSPACE-complete}

Schild says that Proposition~2 follows immediately from Theorem~1, Lemma~1,
Ladner's PSPACE-completeness result for one-modal \(\mathbf K\) validity, and
the extension to every \(m\geq 1\) noted by Halpern and Moses.  The article
defines PSPACE-completeness using log-space many-one reductions on p.~466 and
does not print separate upper- and lower-bound proofs.  The following
two-direction explanation is the handbook's reconstruction.  For the upper
bound, translate \(C\sqsubseteq D\) into validity of
\(f(C)\rightarrow f(D)\); this constructor-wise translation is computable in
logarithmic workspace, and normal multimodal \(\mathbf K_{(m)}\) validity is in
PSPACE.  For the lower bound, one role is already enough: an arbitrary
one-modal formula \(\alpha\) has an inverse concept description \(C_\alpha\),
and \(\alpha\) is valid exactly when
\(\mathsf{Top}\sqsubseteq C_\alpha\).  The inverse constructor-wise translation
is likewise log-space computable, so the one-modal lower bound transfers under
Schild's stated convention without adding roles.

The literal statement of Proposition~2 is that deciding subsumption in
\(\mathcal{ALC}\) is PSPACE-complete even when the vocabulary contains only
one atomic role, under the paper's log-space many-one reduction convention
\cite[pp.~466, 468]{Schild1991}.  In the sentence following the proposition,
Schild separately reports, by referral to Ladner, a
deterministic quadratic-space algorithm for \(\mathcal{ALC}\) coherence.
Lemma~1's logspace, linear-output reductions, together with exchanging the
accept and reject outcomes of a deterministic coherence decider when
incoherence is required, transfer the same quadratic-space upper bound to
equivalence and subsumption.

This presentation concerns only the displayed \(\mathcal{ALC}\) language and
its correspondence with models in which every atomic role is interpreted as
an arbitrary binary relation.  It contains no TBox, ABox, nominal, inverse
role, number restriction, role composition, or global frame constraint.
Neither this \(\mathcal{ALC}\) language nor Schild's Theorem~1 is an OWL
language or OWL profile, and the 1991 paper makes no OWL identity claim.

Those exclusions delimit this presentation, not the whole six-page article.
In Sections~4--4.2 Schild separately treats role composition, union,
reflexive-transitive closure, inverse, identity and range restriction,
features, and concept equations and inequations, with further results about
the resulting terminological logics.  Those later constructions require their
own syntax, semantics, and proof ownership; none is obtained merely by reusing
the displayed structural-induction equivalence
\cite[\S\S4--4.2, pp.~468--471]{Schild1991}.
